% ═══════════════════════════════════════════════════════════════
% 民事起诉状 · 共享说明框
% 通过 % ═══════════════════════════════════════════════════════════════
% 民事起诉状 · 共享说明框
% 通过 % ═══════════════════════════════════════════════════════════════
% 民事起诉状 · 共享说明框
% 通过 % ═══════════════════════════════════════════════════════════════
% 民事起诉状 · 共享说明框
% 通过 \input{../common/notesection} 引入
% ═══════════════════════════════════════════════════════════════

\begin{notesection}
\hfuzz=50pt
\color{inkblack}\heiti\fontsize{10.5bp}{16bp}\selectfont 说明：\par
\vspace*{-0.2bp}
\songti\mdseries\fontsize{10.5bp}{16bp}\selectfont
\setlength{\parindent}{20.7bp}
\sloppy

\setlength{\parskip}{0pt}
\vspace*{-5.1bp}为了方便您更好地参加诉讼，保护您的合法权利，请填写本表。
\lawpar
1.起诉时需向人民法院提交证明您身份的材料，如身份证复印件、营业执照复印件等。
\lawpar
2.本表所列内容是您提起诉讼以及人民法院查明案件事实所需，请务必如实填写。
\lawpar
3.本表有些内容可能与您的案件无关，您认为与案件无关的项目可以填"无"或不填；对于本\\[-4.86bp]
表中勾选项可以在对应项打"√"；您认为另有重要内容需要列明的，可以另附页填写。
\lawpar
4.本表 word 电子版填写时，相关栏目可复制粘贴或扩容，但不得改变要素内容、格式设置。例\\[-4.86bp]
如，多原告、多被告或多委托诉讼代理人等情况，可根据实际情况复制粘贴；需填写文字较多时，\\[-4.86bp]
可根据实际对栏目进行扩容等。
\lawpar
{\fzssk ★}特别提示{\fzssk ★}\\[-4.86bp]
\hspace*{20.7bp}诉讼参加人应遵守诚信原则如实认真填写表格。\\[-4.86bp]
\hspace*{20.7bp}如果诉讼参加人违反有关规定，虚假诉讼、恶意诉讼、滥用诉权，人民法院将视违法情形依法\\[-4.86bp]
追究责任。
\end{notesection}
 引入
% ═══════════════════════════════════════════════════════════════

\begin{notesection}
\hfuzz=50pt
\color{inkblack}\heiti\fontsize{10.5bp}{16bp}\selectfont 说明：\par
\vspace*{-0.2bp}
\songti\mdseries\fontsize{10.5bp}{16bp}\selectfont
\setlength{\parindent}{20.7bp}
\sloppy

\setlength{\parskip}{0pt}
\vspace*{-5.1bp}为了方便您更好地参加诉讼，保护您的合法权利，请填写本表。
\lawpar
1.起诉时需向人民法院提交证明您身份的材料，如身份证复印件、营业执照复印件等。
\lawpar
2.本表所列内容是您提起诉讼以及人民法院查明案件事实所需，请务必如实填写。
\lawpar
3.本表有些内容可能与您的案件无关，您认为与案件无关的项目可以填"无"或不填；对于本\\[-4.86bp]
表中勾选项可以在对应项打"√"；您认为另有重要内容需要列明的，可以另附页填写。
\lawpar
4.本表 word 电子版填写时，相关栏目可复制粘贴或扩容，但不得改变要素内容、格式设置。例\\[-4.86bp]
如，多原告、多被告或多委托诉讼代理人等情况，可根据实际情况复制粘贴；需填写文字较多时，\\[-4.86bp]
可根据实际对栏目进行扩容等。
\lawpar
{\fzssk ★}特别提示{\fzssk ★}\\[-4.86bp]
\hspace*{20.7bp}诉讼参加人应遵守诚信原则如实认真填写表格。\\[-4.86bp]
\hspace*{20.7bp}如果诉讼参加人违反有关规定，虚假诉讼、恶意诉讼、滥用诉权，人民法院将视违法情形依法\\[-4.86bp]
追究责任。
\end{notesection}
 引入
% ═══════════════════════════════════════════════════════════════

\begin{notesection}
\hfuzz=50pt
\color{inkblack}\heiti\fontsize{10.5bp}{16bp}\selectfont 说明：\par
\vspace*{-0.2bp}
\songti\mdseries\fontsize{10.5bp}{16bp}\selectfont
\setlength{\parindent}{20.7bp}
\sloppy

\setlength{\parskip}{0pt}
\vspace*{-5.1bp}为了方便您更好地参加诉讼，保护您的合法权利，请填写本表。
\lawpar
1.起诉时需向人民法院提交证明您身份的材料，如身份证复印件、营业执照复印件等。
\lawpar
2.本表所列内容是您提起诉讼以及人民法院查明案件事实所需，请务必如实填写。
\lawpar
3.本表有些内容可能与您的案件无关，您认为与案件无关的项目可以填"无"或不填；对于本\\[-4.86bp]
表中勾选项可以在对应项打"√"；您认为另有重要内容需要列明的，可以另附页填写。
\lawpar
4.本表 word 电子版填写时，相关栏目可复制粘贴或扩容，但不得改变要素内容、格式设置。例\\[-4.86bp]
如，多原告、多被告或多委托诉讼代理人等情况，可根据实际情况复制粘贴；需填写文字较多时，\\[-4.86bp]
可根据实际对栏目进行扩容等。
\lawpar
{\fzssk ★}特别提示{\fzssk ★}\\[-4.86bp]
\hspace*{20.7bp}诉讼参加人应遵守诚信原则如实认真填写表格。\\[-4.86bp]
\hspace*{20.7bp}如果诉讼参加人违反有关规定，虚假诉讼、恶意诉讼、滥用诉权，人民法院将视违法情形依法\\[-4.86bp]
追究责任。
\end{notesection}
 引入
% ═══════════════════════════════════════════════════════════════

\begin{notesection}
\hfuzz=50pt
\color{inkblack}\heiti\fontsize{10.5bp}{16bp}\selectfont 说明：\par
\vspace*{-0.2bp}
\songti\mdseries\fontsize{10.5bp}{16bp}\selectfont
\setlength{\parindent}{20.7bp}
\sloppy

\setlength{\parskip}{0pt}
\vspace*{-5.1bp}为了方便您更好地参加诉讼，保护您的合法权利，请填写本表。
\lawpar
1.起诉时需向人民法院提交证明您身份的材料，如身份证复印件、营业执照复印件等。
\lawpar
2.本表所列内容是您提起诉讼以及人民法院查明案件事实所需，请务必如实填写。
\lawpar
3.本表有些内容可能与您的案件无关，您认为与案件无关的项目可以填"无"或不填；对于本\\[-4.86bp]
表中勾选项可以在对应项打"√"；您认为另有重要内容需要列明的，可以另附页填写。
\lawpar
4.本表 word 电子版填写时，相关栏目可复制粘贴或扩容，但不得改变要素内容、格式设置。例\\[-4.86bp]
如，多原告、多被告或多委托诉讼代理人等情况，可根据实际情况复制粘贴；需填写文字较多时，\\[-4.86bp]
可根据实际对栏目进行扩容等。
\lawpar
{\fzssk ★}特别提示{\fzssk ★}\\[-4.86bp]
\hspace*{20.7bp}诉讼参加人应遵守诚信原则如实认真填写表格。\\[-4.86bp]
\hspace*{20.7bp}如果诉讼参加人违反有关规定，虚假诉讼、恶意诉讼、滥用诉权，人民法院将视违法情形依法\\[-4.86bp]
追究责任。
\end{notesection}
